%% kapa: exam=1 round=34 year=2023 period=1 form=A subject=civil_law subject_ko=민법 seq=13 number=13
%% source: Q-Net 감정평가사 기출문제 (cst00309) artlSeq=5211789 file=raw/1008/5211789/2023년도 제34회 감정평가사 1차시험 1교시.pdf
%% conversion: kordoc 4.24.51 -> tools/gen1.py (auto); review=proofread-sonnet; flags=none
\begin{question}{13}
의사표시의 효력발생에 관한 설명으로 옳지 않은 것은? (다툼이 있으면 판례에 따름)
\begin{choices}
  \item 의사표시의 발신 후 표의자가 사망하였다면, 그 의사표시는 상대방에게 도달하더라도 무효이다.
  \item 의사표시의 효력발생시기에 관해 도달주의를 규정하고 있는 민법 제111조는 임의규정이다.
  \item 상대방이 정당한 사유 없이 의사표시의 수령을 거절하더라도 상대방이 그 의사표시의 내용을 알 수 있는 객관적 상태에 놓여 있다면 그 의사표시는 효력이 있다.
  \item 재단법인 설립행위의 효력발생을 위해서는 의사표시의 도달이 요구되지 않는다.
  \item 미성년자는 그 행위능력이 제한되고 있는 범위에서 수령무능력자이다.
\end{choices}
\end{question}
\answer{1}
